% common/pedagogy.tex
% Single definition of the shared teaching devices used across all lectures.
% Input this ONCE per lecture, after the theme, e.g.:
%     \usepackage{theme/beamerthemealgorithms}
%     % common/pedagogy.tex
% Single definition of the shared teaching devices used across all lectures.
% Input this ONCE per lecture, after the theme, e.g.:
%     \usepackage{theme/beamerthemealgorithms}
%     % common/pedagogy.tex
% Single definition of the shared teaching devices used across all lectures.
% Input this ONCE per lecture, after the theme, e.g.:
%     \usepackage{theme/beamerthemealgorithms}
%     % common/pedagogy.tex
% Single definition of the shared teaching devices used across all lectures.
% Input this ONCE per lecture, after the theme, e.g.:
%     \usepackage{theme/beamerthemealgorithms}
%     \input{common/pedagogy.tex}
% The lecture-specific *TikZ styles* stay in each lectureNN/common.tex;
% only the pedagogy wrappers below are centralized here.

% ---- Canonical devices -------------------------------------------------
\newcommand{\checkpoint}[1]{\begin{exampleblock}{Checkpoint}#1\end{exampleblock}}
\newcommand{\mission}[1]{\begin{block}{Mission}#1\end{block}}
% \takeaway environment and \keyterm/\muted come from the theme.

% ---- Pseudocode size floor ---------------------------------------------
% One knob for every algorithmic block. Never smaller than \small.
% Usage:  \begin{frame}{...}\algofont \begin{algorithmic}[1] ... \end{algorithmic}\end{frame}
\newcommand{\algofont}{\small}
% Optional: wrap a whole block so the size is impossible to forget.
\newenvironment{algoblock}[1][1]{\par\algofont\begin{algorithmic}[#1]}{\end{algorithmic}}

% ---- "Optional / 선택 심화" badge (was defined only in lecture02) --------
\providecommand{\optional}[1][Optional]{%
  \tikz[baseline=(label.base)]\node[draw=AlgoRule,rounded corners=2pt,
    fill=AlgoLight,text=AlgoGrayText,font=\scriptsize\bfseries,
    inner xsep=4pt,inner ysep=2pt] (label) {#1};}

% ---- Backward-compatible aliases ---------------------------------------
% So existing section files keep compiling unchanged. Remove these once the
% section files are migrated to the canonical names above.
\newcommand{\sectiontakeaway}[1]{\begin{takeaway}#1\end{takeaway}}
\newcommand{\sortingtakeaway}[1]{\begin{takeaway}#1\end{takeaway}}
\newcommand{\selecttakeaway}[1]{\begin{takeaway}#1\end{takeaway}}
\newcommand{\dptakeaway}[1]{\begin{takeaway}#1\end{takeaway}}
\newcommand{\sttakeaway}[1]{\begin{takeaway}#1\end{takeaway}}
\newcommand{\httakeaway}[1]{\begin{takeaway}#1\end{takeaway}}
\newcommand{\gtakeaway}[1]{\begin{takeaway}#1\end{takeaway}}
\newcommand{\stakeaway}[1]{\begin{takeaway}#1\end{takeaway}}
\newcommand{\sstakeaway}[1]{\begin{takeaway}#1\end{takeaway}}
\newcommand{\dpcheck}[1]{\checkpoint{#1}}
\newcommand{\htcheckpoint}[1]{\checkpoint{#1}}
\newcommand{\stcheckpoint}[1]{\checkpoint{#1}}
\newcommand{\gcheckpoint}[1]{\checkpoint{#1}}
\newcommand{\scheckpoint}[1]{\checkpoint{#1}}
\newcommand{\sscheckpoint}[1]{\checkpoint{#1}}
\newcommand{\htmission}[1]{\mission{#1}}
\newcommand{\gmission}[1]{\mission{#1}}
\newcommand{\smission}[1]{\mission{#1}}
\newcommand{\ssmission}[1]{\mission{#1}}
\newcommand{\ssoptional}[1][Optional]{\optional[#1]}

% The lecture-specific *TikZ styles* stay in each lectureNN/common.tex;
% only the pedagogy wrappers below are centralized here.

% ---- Canonical devices -------------------------------------------------
\newcommand{\checkpoint}[1]{\begin{exampleblock}{Checkpoint}#1\end{exampleblock}}
\newcommand{\mission}[1]{\begin{block}{Mission}#1\end{block}}
% \takeaway environment and \keyterm/\muted come from the theme.

% ---- Pseudocode size floor ---------------------------------------------
% One knob for every algorithmic block. Never smaller than \small.
% Usage:  \begin{frame}{...}\algofont \begin{algorithmic}[1] ... \end{algorithmic}\end{frame}
\newcommand{\algofont}{\small}
% Optional: wrap a whole block so the size is impossible to forget.
\newenvironment{algoblock}[1][1]{\par\algofont\begin{algorithmic}[#1]}{\end{algorithmic}}

% ---- "Optional / 선택 심화" badge (was defined only in lecture02) --------
\providecommand{\optional}[1][Optional]{%
  \tikz[baseline=(label.base)]\node[draw=AlgoRule,rounded corners=2pt,
    fill=AlgoLight,text=AlgoGrayText,font=\scriptsize\bfseries,
    inner xsep=4pt,inner ysep=2pt] (label) {#1};}

% ---- Backward-compatible aliases ---------------------------------------
% So existing section files keep compiling unchanged. Remove these once the
% section files are migrated to the canonical names above.
\newcommand{\sectiontakeaway}[1]{\begin{takeaway}#1\end{takeaway}}
\newcommand{\sortingtakeaway}[1]{\begin{takeaway}#1\end{takeaway}}
\newcommand{\selecttakeaway}[1]{\begin{takeaway}#1\end{takeaway}}
\newcommand{\dptakeaway}[1]{\begin{takeaway}#1\end{takeaway}}
\newcommand{\sttakeaway}[1]{\begin{takeaway}#1\end{takeaway}}
\newcommand{\httakeaway}[1]{\begin{takeaway}#1\end{takeaway}}
\newcommand{\gtakeaway}[1]{\begin{takeaway}#1\end{takeaway}}
\newcommand{\stakeaway}[1]{\begin{takeaway}#1\end{takeaway}}
\newcommand{\sstakeaway}[1]{\begin{takeaway}#1\end{takeaway}}
\newcommand{\dpcheck}[1]{\checkpoint{#1}}
\newcommand{\htcheckpoint}[1]{\checkpoint{#1}}
\newcommand{\stcheckpoint}[1]{\checkpoint{#1}}
\newcommand{\gcheckpoint}[1]{\checkpoint{#1}}
\newcommand{\scheckpoint}[1]{\checkpoint{#1}}
\newcommand{\sscheckpoint}[1]{\checkpoint{#1}}
\newcommand{\htmission}[1]{\mission{#1}}
\newcommand{\gmission}[1]{\mission{#1}}
\newcommand{\smission}[1]{\mission{#1}}
\newcommand{\ssmission}[1]{\mission{#1}}
\newcommand{\ssoptional}[1][Optional]{\optional[#1]}

% The lecture-specific *TikZ styles* stay in each lectureNN/common.tex;
% only the pedagogy wrappers below are centralized here.

% ---- Canonical devices -------------------------------------------------
\newcommand{\checkpoint}[1]{\begin{exampleblock}{Checkpoint}#1\end{exampleblock}}
\newcommand{\mission}[1]{\begin{block}{Mission}#1\end{block}}
% \takeaway environment and \keyterm/\muted come from the theme.

% ---- Pseudocode size floor ---------------------------------------------
% One knob for every algorithmic block. Never smaller than \small.
% Usage:  \begin{frame}{...}\algofont \begin{algorithmic}[1] ... \end{algorithmic}\end{frame}
\newcommand{\algofont}{\small}
% Optional: wrap a whole block so the size is impossible to forget.
\newenvironment{algoblock}[1][1]{\par\algofont\begin{algorithmic}[#1]}{\end{algorithmic}}

% ---- "Optional / 선택 심화" badge (was defined only in lecture02) --------
\providecommand{\optional}[1][Optional]{%
  \tikz[baseline=(label.base)]\node[draw=AlgoRule,rounded corners=2pt,
    fill=AlgoLight,text=AlgoGrayText,font=\scriptsize\bfseries,
    inner xsep=4pt,inner ysep=2pt] (label) {#1};}

% ---- Backward-compatible aliases ---------------------------------------
% So existing section files keep compiling unchanged. Remove these once the
% section files are migrated to the canonical names above.
\newcommand{\sectiontakeaway}[1]{\begin{takeaway}#1\end{takeaway}}
\newcommand{\sortingtakeaway}[1]{\begin{takeaway}#1\end{takeaway}}
\newcommand{\selecttakeaway}[1]{\begin{takeaway}#1\end{takeaway}}
\newcommand{\dptakeaway}[1]{\begin{takeaway}#1\end{takeaway}}
\newcommand{\sttakeaway}[1]{\begin{takeaway}#1\end{takeaway}}
\newcommand{\httakeaway}[1]{\begin{takeaway}#1\end{takeaway}}
\newcommand{\gtakeaway}[1]{\begin{takeaway}#1\end{takeaway}}
\newcommand{\stakeaway}[1]{\begin{takeaway}#1\end{takeaway}}
\newcommand{\sstakeaway}[1]{\begin{takeaway}#1\end{takeaway}}
\newcommand{\dpcheck}[1]{\checkpoint{#1}}
\newcommand{\htcheckpoint}[1]{\checkpoint{#1}}
\newcommand{\stcheckpoint}[1]{\checkpoint{#1}}
\newcommand{\gcheckpoint}[1]{\checkpoint{#1}}
\newcommand{\scheckpoint}[1]{\checkpoint{#1}}
\newcommand{\sscheckpoint}[1]{\checkpoint{#1}}
\newcommand{\htmission}[1]{\mission{#1}}
\newcommand{\gmission}[1]{\mission{#1}}
\newcommand{\smission}[1]{\mission{#1}}
\newcommand{\ssmission}[1]{\mission{#1}}
\newcommand{\ssoptional}[1][Optional]{\optional[#1]}

% The lecture-specific *TikZ styles* stay in each lectureNN/common.tex;
% only the pedagogy wrappers below are centralized here.

% ---- Canonical devices -------------------------------------------------
\newcommand{\checkpoint}[1]{\begin{exampleblock}{Checkpoint}#1\end{exampleblock}}
\newcommand{\mission}[1]{\begin{block}{Mission}#1\end{block}}
% \takeaway environment and \keyterm/\muted come from the theme.

% ---- Pseudocode size floor ---------------------------------------------
% One knob for every algorithmic block. Never smaller than \small.
% Usage:  \begin{frame}{...}\algofont \begin{algorithmic}[1] ... \end{algorithmic}\end{frame}
\newcommand{\algofont}{\small}
% Optional: wrap a whole block so the size is impossible to forget.
\newenvironment{algoblock}[1][1]{\par\algofont\begin{algorithmic}[#1]}{\end{algorithmic}}

% ---- "Optional / 선택 심화" badge (was defined only in lecture02) --------
\providecommand{\optional}[1][Optional]{%
  \tikz[baseline=(label.base)]\node[draw=AlgoRule,rounded corners=2pt,
    fill=AlgoLight,text=AlgoGrayText,font=\scriptsize\bfseries,
    inner xsep=4pt,inner ysep=2pt] (label) {#1};}

% ---- Backward-compatible aliases ---------------------------------------
% So existing section files keep compiling unchanged. Remove these once the
% section files are migrated to the canonical names above.
\newcommand{\sectiontakeaway}[1]{\begin{takeaway}#1\end{takeaway}}
\newcommand{\sortingtakeaway}[1]{\begin{takeaway}#1\end{takeaway}}
\newcommand{\selecttakeaway}[1]{\begin{takeaway}#1\end{takeaway}}
\newcommand{\dptakeaway}[1]{\begin{takeaway}#1\end{takeaway}}
\newcommand{\sttakeaway}[1]{\begin{takeaway}#1\end{takeaway}}
\newcommand{\httakeaway}[1]{\begin{takeaway}#1\end{takeaway}}
\newcommand{\gtakeaway}[1]{\begin{takeaway}#1\end{takeaway}}
\newcommand{\stakeaway}[1]{\begin{takeaway}#1\end{takeaway}}
\newcommand{\sstakeaway}[1]{\begin{takeaway}#1\end{takeaway}}
\newcommand{\dpcheck}[1]{\checkpoint{#1}}
\newcommand{\htcheckpoint}[1]{\checkpoint{#1}}
\newcommand{\stcheckpoint}[1]{\checkpoint{#1}}
\newcommand{\gcheckpoint}[1]{\checkpoint{#1}}
\newcommand{\scheckpoint}[1]{\checkpoint{#1}}
\newcommand{\sscheckpoint}[1]{\checkpoint{#1}}
\newcommand{\htmission}[1]{\mission{#1}}
\newcommand{\gmission}[1]{\mission{#1}}
\newcommand{\smission}[1]{\mission{#1}}
\newcommand{\ssmission}[1]{\mission{#1}}
\newcommand{\ssoptional}[1][Optional]{\optional[#1]}
